\documentclass{handout}

\assignment{Lab 2}
\title{Basics of motion}

\begin{document}
\maketitle

The study of motion is a study of how external influences (forces) affect an object's motion. The starting point of this analysis, however, is establishing a baseline: what does the motion of an object look like in the absence of any interactions? Here are two theories, both of which play an important role in the history of physics:
\begin{enumerate}
    \item An object's natural tendency, in the absence of interactions with other objects, is to \textbf{maintain its current motion in a straight line}. If it is not moving, it will remain not moving. If it is moving in a straight line, it will go in that straight line forever at the same speed.
    \item An object's natural tendency, in the absence of interactions with other objects, \textbf{is to be stationary}. Any motion is because of some interaction causing it to maintain it motion (e.g.~the air rushing back into the space pushes the object from behind or internal interactions like an organism's  metabolism), but this effect weakens eventually and the object returns to rest.
\end{enumerate}

Today, we would like to try to adjudicate between these two possibilities. There is one
problem: \textbf{when was the last time you were able to observe an
object that was not interacting with any other objects?}

\section{\textit{Gedankenexperiment} of non-interaction}\label{the-limit-of-no-interaction}

\emph{Gedankenexperiment} is a German word that translates to `thought experiment.' For most of the history of physics, this was the primary way of doing physics. The real world is messy: there are a lot of complex systems of cause and effect operating simultaneously. In a thought experiment, we \emph{imagine} a scenario that is as simple as possible (in fact, usually so simple as to be \emph{impossible}!).

Usually, the impossibility of a thought experiment occurs in applying the concept of a \textbf{limit}. Sometimes we can see a trend: e.g. ``what happen as an interaction becomes weaker?'' and make a limiting statement: e.g. ``what would happen if the interaction were not there at all?'' The latter does not need to be possible in the real world. Instead, it functions as a tool to help us think about physical relationships, serving as starting points for physical laws.

Consider running to a constant speed of 3 m/s and sliding down the hallway in each of the following scenarios:
\begin{enumerate}
    \item Sliding across the floor on a hovercraft
    \item Sliding across the floor with shoes on
    \item Rolling across the floor on an office chair
    \item Sliding across the floor with socks on
\end{enumerate}

\subsection{Discussion questions}
\begin{questions}
    \question Rank these four scenarios in terms of how far you slide

    \ranking[4]{Farthest}{Shortest}

    \question Rank these four combinations in terms of how strong the interaction is
between you and the floor

    \ranking[4]{Strongest}{Weakest}


    \question Is your interaction with the air different in these scenarios? Why or why
not? Is this effect significant?

    \answerspace{1.5in}

    \question Imagine a hypothetical super puck sliding on a surface. This super puck
does not interact with anything. The air does not interact with the puck
and the puck does not interact with the surface \textbf{at all}. Use
your answers above to explain what you believe the motion of this
unreal super puck \emph{would} be like. Justify your answer using a limit.

    \answerspace{1.5in}

    \question Which of the 2 theories of motion described in section 1 do you think is correct and why?

    \answerspace{1.5in}

\end{questions}

\section{Force and frictionless
motion}\label{force-and-frictionless-motion}

Now we will be using the hovercrafts. Operation is simple
\begin{itemize}
    \item  Sit in the middle of the hovercraft.
\item Turn on the blower.
\item Have a friend push you to get you moving.
\item Make sure there is someone at the end of the
hallway to stop you.
\end{itemize}
Try a couple people of different size, try speeding them up, stopping them, and turning the hovercraft. Think about the forces that cause all these, direction, magnitude, etc.


\subsection{Discussion questions}
\begin{questions}
\question Does your experience of hovercraft motion match your identified theory of motion from part I?

\answerspace{1in}

\question We would like to determine what the effects of interaction could be. Try
to figure out qualitatively what the effect of the following
interactions are:
\begin{parts}
    \part A force pointing in the same direction the hovercraft is moving.

    \answerspace{1in}

    \part A force pointing in the opposite direction the hovercraft is moving.

    \answerspace{1in}

    \part A force perpendicular to the motion of the hovercraft.

    \answerspace{1in}
\end{parts}

\question How is each of these interactions affected by the \textbf{strength} of
the interaction (how hard one pushes)?

\answerspace{1.5in}

\question How are each of these interactions affected by the \textbf{mass} of the
hovercraft operator (how heavy something is)?

\answerspace{1.5in}

\question In what sense was your exploration of the hovercraft an \textbf{experiment} as opposed to a simple \textbf{observation}? 

\answerspace{1.5in}

\end{questions}

\clearpage
\section{Quantifying motion and
forces}\label{quantifying-motion-and-forces}
In our first section, we arrived at Newton's first law of motion (in the absence of interactions, an object at rest remains at rest and a moving object continues at the same speed in a straight line).

Newton's second law relates three quantities: $F$ is the force applied to the object; $m$ is the mass, a number that quantifies how heavy something is in the sense that it sluggishly reacts to forces; and $a$ is acceleration, the rate of change of velocity (speed taking into account direction).

We can write down proprtional relationships in either of the following ways:
\begin{align*}
    A&= k B \\
A &= \frac{k}{B} 
\end{align*}
The first relationship says that $A$ is \textbf{directly proportional} to $B$; if $B$ doubles, $A$ doubles. The second relationship says that $A$ is \textbf{inversely proportional} to $B$; if $B$ doubles, $A$ is cut in half.

Let's start writing an equation that captures the relationship between force, mass, and acceleration. We'll begin by noting the direct and inverse proportionalities we've observed. 

For example, read the first as asking ``If I double force, and expect acceleration to stay constant, would mass have to double or be cut in half?'' and ``If I double the force applied to an object of constant mass, would acceleration double or be cut in half?''
\answerspace{\baselineskip}

\begin{questions}
    \question \textbf{Force} is \blank[2in] proportional to \textbf{mass} and \blank[2in] proportional to \textbf{acceleration}.

\answerspace{\baselineskip}
    \question \textbf{Mass} is \blank[2in] proportional to \textbf{force} and \blank[2in] proportional to \textbf{acceleration}.

\answerspace{\baselineskip}

    \question \textbf{Acceleration} is \blank[2in] proportional to \textbf{mass} and \blank[2in] proportional to \textbf{force}.

\answerspace{\baselineskip}
\clearpage
    \question \label{q:firstlaw} Write a mathematical expression that captures this relationship and explain how you chose this form. 

    \answerspace{\fill}

    \question Explain how negative numbers help capture the distinction between starting and stopping.

    \answerspace{\fill}

    \question Explain how this equation does or does not capture the notion of turning. 

    \answerspace{\fill}
    \answerspace{.5in}

    \question In order to properly treat all possible applications of the equation from question \ref{q:firstlaw}, we would need an algebraic language capable of describing simultaneously\\~ \\ \blank[2in] and \blank[2in].


\end{questions}

\end{document}
